\honest{Einstein's Superpowers}{Eliezer Yudkowsky}{2008}

There is a widespread tendency to talk as if Einstein and Newton had superpowers. I used to think so too, until Julian Barbour's \textsc{The End of Time} cured me. Between Barbour's polite lines I thought I heard him shouting: ``What Einstein did isn't magic, people!'' \nb{In a 2013 edit Yudkowsky added that Barbour did not actually say this; the post had already offered it as the author's own reading.} I kinda suspect that people once sniffed at Barbour, ``Oh, you think you're Einstein, do you?'' I give no evidence for this.

John Baez's Crackpot Index gives points for comparing yourself to Einstein. I cite it to show ``what society generally thinks'' of such comparisons. \nb{The Index is one physicist's joke list. It shows Baez's experience, not what society thinks.}

The crackpot thinks Einstein had superpowers and claims them too. The critic who is shocked when a young physicist hopes to do better than Einstein makes the same mistake from the other side. Both confuse social status with research potential, and imagine that Einstein's potential was as rare as his fame.

Not everyone can be Einstein, I grant, ``in the modern world.'' But bright people are not all that rare, and many end up ``just another Jewish genius.'' What separates them from Einstein, I suspect, is not brilliance but failing to choose an important problem, to find a worthwhile angle of attack, and to persist for years without support. \nb{Richard Hamming had said much the same in ``You and Your Research'' (1986): great work ``is something else than mere brains,'' and ``If you do not work on an important problem, it's unlikely you'll do important work.'' The post does not mention him.} Had you met Albert before his papers, you would have seen no aura of destiny. \nb{The record agrees: almost two years without a teaching post, then an assistant examiner's job at the Swiss Patent Office.}

All this matters to me because I have a genuinely important problem, an angle of attack, years of work behind me, and a support structure, and people keep saying, ``Let's see your aura of destiny, buddy.''

Barbour saw through Einstein, and said ``something along the lines of'' where Einstein ``failed to apply his own methods, and missed the key insight.'' Einstein was no ordinary bloke. I take a special joy in doing things people call ``humanly impossible.'' \nb{The link is to the author's AI-Box Experiment, which its own page calls ``strictly anecdotal evidence.'' In a later post Yudkowsky reported three more games: ``I won the first, and then lost the next two.''} You acquire such powers by seeing that they are perfectly normal. ``This is a general principle in life.''
